\chapter{Requirements Analysis}
\label{ch:requirements}

\section{Stakeholders and Roles}
\label{sec:stakeholders}

\thesisfigure{f4-use-case}
  {Actors and goals as implemented. Shaded use cases are unreachable in a
   default deployment because every feature flag defaults to off.}
  {fig:use-case}

\TODO{Two things in \cref{fig:use-case} deserve prose rather than being left to
the diagram. First, content authoring is gated on \emph{ownership}, not role:
\code{require\_creator} admits both professors and students, and the code
comment states the intent explicitly. Second, that decision has a direct
security consequence --- a student-uploaded PDF can reach another student's
tutor context --- which is why retrieved text is treated as untrusted in
\cref{sec:injection}. Requirements decisions with downstream security
consequences are exactly what an examiner looks for.}

\section{Functional Requirements}
\label{sec:functional-requirements}

\TODO{Derive from the use cases, numbered FR1..FRn so \cref{ch:evaluation} can
reference them. Keep each testable --- "the system shall reject an upload
exceeding the configured page limit" is testable; "the system shall be
user-friendly" is not.}

\section{Non-Functional Requirements}
\label{sec:nonfunctional-requirements}

\TODO{The ones this system actually takes a position on, with the real values
from \cref{app:constants}: upload limits, page limits, queue backpressure
threshold, per-user LLM cost cap, and the resilience target for provider
failure. Numbered NFR1..NFRn.}

\TODO{Include idempotency as an explicit non-functional requirement. It is
unusual for a bachelor thesis to name it, and it is the property
\cref{sec:idempotency} shows is satisfied by six independent mechanisms --- so
stating the requirement here makes that section an answer rather than a
digression.}

\section{Domain Model}
\label{sec:domain-model}

\thesisfigure{f5-domain-model}
  {Conceptual domain model. Analysis-level vocabulary, deliberately distinct
   from the physical schema in \cref{fig:er}.}
  {fig:domain-model}

\TODO{Note explicitly that \emph{Quiz} appears in \cref{fig:domain-model} as
domain vocabulary but has no table behind it --- a quiz is the set of questions
reachable through a lecture's slides. Naming a concept that deliberately has no
persistent representation shows the distinction between analysis and
implementation model is understood, rather than accidental.}

\section{Constraints}
\label{sec:constraints}

\TODO{The real ones, because they shaped the design and make the trade-offs in
\cref{ch:design} legible: a single small virtual machine (4\,GB RAM, 2 vCPU),
free-tier or low-cost LLM provider quotas, a managed Postgres instance with no
backups on the free tier, and a one-person development team.}
